% =====================================================================
% BAB V KESIMPULAN DAN SARAN
% =====================================================================

\chapter{KESIMPULAN DAN SARAN}
\label{ch:kesimpulan}

\section{Kesimpulan}
\label{sec:kesimpulan}

Tugas akhir ini mengembangkan AXIS sebagai \textit{companionship chatbot}
non-klinis untuk mahasiswa Indonesia. Pengembangan diarahkan oleh tiga
rumusan masalah: membangun percakapan pendamping yang empatik dan aman,
mengorkestrasi mood harian serta asesmen skrining depresi PHQ-9 melalui
percakapan, dan menjaga kesinambungan konteks lintas sesi melalui memori
jangka panjang berbasis \textit{knowledge graph}. Berdasarkan rancangan,
implementasi, dan evaluasi pada Bab IV, diperoleh kesimpulan berikut.

Pertama, AXIS menggabungkan percakapan reflektif berprinsip CBT, pengayaan
ragam bahasa, dan batas keselamatan dalam satu alur percakapan. Validasi
kontrak menunjukkan 54 skenario kebijakan dialog CBT dan 80 skenario
\textit{guardrail} lulus. Pada penilaian buta berskala kecil, AXIS dipilih pada
10 dari 18 penilaian berpasangan; perbedaan yang paling konsisten terletak pada
ketepatan penggunaan konteks memori, bukan pada seluruh dimensi kualitas
percakapan. Pada \textit{benchmark} 50 pesan keselamatan, sistem mencapai
sensitivitas 0,654 dan spesifisitas 0,917. Seluruh respons pada kasus yang
berhasil dideteksi memenuhi rubrik keselamatan, tetapi sembilan pesan berisiko
masih tidak terdeteksi. Oleh sebab itu, hasil ini membuktikan bahwa mekanisme
keselamatan dan percakapan telah berjalan pada skenario yang diuji, bukan bahwa
sistem sudah memahami seluruh variasi ekspresi krisis atau lebih empatik secara
umum.

Kedua, mood harian dan PHQ-9 telah diorkestrasi sebagai bagian dari
percakapan, bukan sebagai diagnosis maupun formulir yang terlepas dari konteks
pengguna. Sebanyak 130 kontrak otomatis yang mencakup penawaran, penolakan,
penyampaian item, klarifikasi, penskoran, umpan balik, dan rute item kesembilan
lulus. Pada korpus 80 jawaban bebas, 73 jawaban dengan frekuensi tegas dipetakan
sesuai label acuan dengan \textit{macro-F1} dan \textit{quadratic weighted
kappa} sebesar 1,00, sedangkan tujuh jawaban ambigu diarahkan ke klarifikasi.
Hasil ini mendukung kesetiaan orkestrasi dan pemetaan jawaban bebas pada korpus
uji, tetapi belum membuktikan validitas psikometrik maupun manfaat klinis dari
sistem.

Ketiga, arsitektur memori hibrid memungkinkan sistem menyimpan dan membaca
kembali konteks secara relasional. Sebanyak 47 kontrak memori inti dan 27
kontrak \textit{lifecycle} serta konteks lulus pada skenario deterministik.
Pada 15 kueri parafrase satu-\textit{hop}, kondisi hibrid dan vektor-saja
memperoleh \textit{Recall@5} yang sama, yaitu 1,00. Sebaliknya, pada satu kueri
yang membutuhkan identitas relasi graf, kondisi hibrid berhasil menemukan
pengalaman terkait sedangkan kondisi vektor-saja tidak. Probe pemaknaan ulang
ujung-ke-ujung juga menunjukkan satu dari lima relasi yang ditulis selaras
secara semantik menurut satu penilai LLM, dan tiga pikiran lama yang tersupersede
seluruhnya tidak lagi aktif. Temuan ini menunjukkan nilai graf pada kueri
relasional dan pengelolaan siklus hidup memori, tetapi belum cukup untuk
mengklaim keunggulan umum terhadap pencarian semantik atau ketepatan semantik
penulisan memori pada penggunaan nyata.

Secara keseluruhan, AXIS telah menghasilkan purwarupa yang memadukan
pendampingan percakapan, asesmen skrining depresi secara percakapan, dan memori
jangka panjang dalam satu sistem yang dapat diuji. Pengujian pengguna nyata
belum dilakukan; karena itu, kesimpulan mengenai rasa didampingi, kemudahan
penggunaan, dan kenyamanan pengguna belum dapat ditarik pada laporan ini.

\section{Keterbatasan}
\label{sec:keterbatasan}

Evaluasi pada laporan ini masih memiliki beberapa keterbatasan. Pertama,
validasi kontrak membuktikan perilaku sistem pada skenario terkontrol, bukan
seluruh kemungkinan percakapan bebas. Verifikasi teknis untuk unduh data masih
berada pada lapisan antarmuka dan belum memiliki kontrak otomatis tersendiri.

Kedua, penilaian respons terbuka menggunakan satu konfigurasi model sebagai
\textit{LLM-as-judge}. Metode ini membantu menerapkan rubrik secara konsisten,
tetapi tidak menggantikan penilaian pengguna atau ahli. Ukuran korpus dialog,
keselamatan, dan jawaban bebas PHQ-9 juga masih terbatas. Sembilan pesan berisiko
yang luput pada \textit{benchmark} keselamatan menunjukkan bahwa deteksi
eufemisme belum memadai untuk digeneralisasi.

Ketiga, kalibrasi retrieval eksternal sudah menghitung \textit{P@5}, MRR, dan
nDCG@5 pada LongMemEval_S, tetapi korpus tersebut belum menyediakan anotasi
graf AXIS maupun pasangan jawaban berlabel. Oleh sebab itu,
\textit{grounded-answer rate}, \textit{contradiction rate},
\textit{false-personalization rate}, \textit{update correctness}, dan
\textit{stale-memory rate} belum dapat dihitung. Probe pemaknaan ulang juga
hanya terdiri atas lima relasi sehingga tidak dapat menjadi estimasi populasi.

Keempat, pengujian pengguna nyata belum dilaksanakan. Dengan demikian,
pengalaman subjektif mahasiswa, pemahaman terhadap kendali data, dan kemudahan
antarmuka belum memiliki bukti empiris dari partisipan. Sistem ini juga bukan
alat diagnosis dan tidak dievaluasi sebagai intervensi klinis.

\section{Saran}
\label{sec:saran}

Pengembangan berikutnya perlu memprioritaskan pengujian pengguna nyata dengan
instrumen yang telah disiapkan pada Lampiran~\ref{lamp:userstudy}. Pengujian
tersebut perlu mengukur pengalaman pendampingan, kemudahan penggunaan, dan
pemahaman kendali data tanpa memperluasnya menjadi klaim diagnosis.

Korpus evaluasi keselamatan, bahasa informal, bahasa campur, dan eufemisme
perlu diperluas serta ditinjau dengan rubrik yang sama. Pesan risiko yang belum
tertangkap perlu dianalisis untuk memperbaiki aturan dan penilaian risiko.
Penilaian respons juga sebaiknya melibatkan penilai manusia dalam skala yang
memadai agar keselarasan \textit{LLM-as-judge} dapat diperiksa secara langsung.

Pada memori, diperlukan korpus standar emas yang memuat fakta, relasi, konflik,
pembaruan, dan kueri relevansi berjenjang. Korpus tersebut memungkinkan
pengukuran \textit{retrieval} dan \textit{lifecycle} secara lebih lengkap,
termasuk pengujian hibrid dan vektor-saja pada urutan kueri yang identik.
Pemantauan sinkronisasi antara penyimpanan graf dan vektor, serta pengujian
penggunaan jangka panjang, juga perlu dilakukan sebelum sistem digunakan lebih
luas.
